\section{Phase Change：从 loss bump 提出 circuit hypothesis}

二维 loss map 的关键发现不是 ICL 缓慢增强，而是在很窄的训练窗口内突然出现。为了突出这种变化，研究者对 token-position curve 与 training trajectory 做差分/导数；热图中出现清晰竖线，普通 loss curve 也出现一个小 bump。

\lecturefigure{slide-08-per-step-learning-heatmaps.jpg}{Loss surface 与导数视图：窄训练窗口内出现明显结构变化。}{Stanford Online 官方视频 00:13:48--00:14:01。}

可把诊断量抽象写为
\[
D_i(s,i)=\frac{\partial \mathcal{L}(s,i)}{\partial \log i},
\qquad
D_s(s)=\frac{\partial S_{\mathrm{ICL}}(s)}{\partial \log s}.
\]

\begin{itemize}
  \item $D_i$ 衡量增加相对 context 长度带来的 loss reduction；
  \item $D_s$ 衡量 ICL score 在训练中的变化速度；
  \item narrow spike 表明行为在少量训练 steps 内快速改变；
  \item 导数放大噪声，因此必须对照原始 loss、多个 seeds 与多个 model sizes。
\end{itemize}

\lecturefigure{slide-09-phase-change-across-scales.jpg}{不同 model scales 中都能看到 loss bump 与 ICL 快速上升。}{Stanford Online 官方视频 00:14:02--00:15:11。}

\begin{warningbox}{“Phase change” 是经验描述，不是已证明的物理相变}
曲线在有限 snapshots 中快速变化，并不证明 optimization trajectory 数学不连续。更稳妥的结论是：存在一个 narrow training interval，模型行为与内部 signatures 同时快速重组，值得提出机制假说并做干预。
\end{warningbox}

\subsection{从行为变化到内部机制}

如果外部 ICL 在几百 steps 内突然增强，一个自然假设是模型内部也形成了新 circuit。要验证它，不能只继续看 loss；需要把参数分解为可解释 path，并比较一层与两层 attention-only models，因为前者没有该 phase change，后者有。

\lecturefigure{slide-10-sudden-mechanism-change-hypothesis.jpg}{研究假设：ICL 行为相变对应内部算法/电路的突然形成。}{Stanford Online 官方视频 00:15:12--00:15:31。}

\teachervoice{讲者用的是 “natural hypothesis”，不是 proof。行为 bump 只是搜索入口；只有当候选 circuit 在时间上出现、在结构上能实现目标算法、并且 ablation 会移除行为时，机制解释才逐渐成立。}

\subsection{证据阶梯}

后文会反复使用四种证据。为了避免混淆，先给出强度排序：
\begin{table}[H]
\centering
\small
\caption{本讲中的证据类型与可支持结论。}
\begin{tabularx}{\textwidth}{L{0.21\textwidth}L{0.28\textwidth}X}
\toprule
证据 & 例子 & 能支持什么 \\
\midrule
Representational plausibility & QK/OV 矩阵能实现 match 与 copy & 说明 circuit 有能力，不证明实际被使用。 \\
Temporal co-occurrence & Induction signature 与 ICL 同时出现 & 提供相关证据，可能有共同原因。 \\
Behavioral generality & Copy、translation、pattern matching & 说明机制不局限于单一 toy string。 \\
Causal intervention & Ablate head 后 ICL score 回退 & 直接支持该组件对该 metric 的因果贡献。 \\
Near-complete accounting & 少量 paths 解释 toy model 主要 loss & 在限定模型中接近完整，但不自动外推大型 LLM。 \\
\bottomrule
\end{tabularx}
\end{table}

\subsection{本章小结}

Phase change 把一个模糊的“模型似乎会从 context 学习”转成具体研究目标：解释窄窗口内的 ICL score 变化。行为曲线产生 hypothesis，真正的机制结论还需要 algebra、结构可行性、跨例泛化与 causal ablation。

